\chapter{自然语言处理概述}
\label{chap:nlp-introduction}

一份发布于2025年4月15日的活动通知写道：

\begin{quote}
机器学习读书会原定于4月17日举行，现改至4月18日（周五）14:00，地点为海棠楼201。
校内师生可直接参加，校外来宾须在4月16日18:00前提交报名信息。
\end{quote}

这份通知及本章的咨询、数据库、预测分数和执行记录均为教学构造，不是实际活动或实测
结果。活动编号统一为\code{ML-0418}，日期均采用北京时间；2025年4月18日确为周五。
通知正文不含标题、换行和首尾空白，共80个字符，后文的字符位置均以这段正文为准。

面对这段文字，有人想知道活动是否改期，有人需要把时间和地点录入日历，也有人会追问
“我明天下午到可以吗”。系统若只找出“时间”“地点”等关键词，还不能完成这些需求：
它必须知道要交付什么结果，可以参考哪些信息，以及怎样判断结果正确。把自然语言交给
计算机处理，困难不只在于文字形式多样，更在于同一句话会随着任务、上下文和使用目的
承担不同作用。

\term{自然语言处理}（natural language processing, NLP）研究如何用计算方法分析、
组织、生成人类语言，并通过语言完成信息获取与交流任务。这里的“处理”既包括为词语、
句子或文档给出结构化分析，也包括分类、抽取、转换、问答、对话、程序生成和语音交互。
这些任务可以共享模型和训练方法，验收标准却不能由模型名称决定：摘要要保留关键事实，
报名操作要真正改变业务状态，语音输出还要让听者正确听到日期和专名。

本章从上述通知出发，回答三个相连的问题。语言应用为什么需要明确任务，定义任务时要
约定哪些对象与信息条件，不同实现和评价证据又怎样对应这些约定。最后的历史脉络按
研究问题展开，用来说明任务条件如何逐渐扩展；模型结构与学习原理分别回指前面的模型卷
和范式卷。

\section{语言处理的需求与应用场景}
\label{sec:nlp-needs-and-scenarios}

阅读通知时，人的目标会决定关注点。若只想复述通知内容，需要辨认活动、时间、地点和
报名条件；若要维护活动数据库，还要把这些内容写入约定字段；若要回答咨询，则必须结合
咨询者身份、当前日期和已有对话判断哪些事实与问题相关。三种需求分别偏向语言分析、
内容整理与交流，但它们并非彼此隔离。可靠的答复常常要先找到事实，可靠的登记又要保留
事实来自原文的哪一处。

围绕通知，考虑三条具体咨询。甲在4月15日问：“读书会改到什么时候、在哪里？”
合格答复是“2025年4月18日周五14:00，海棠楼201”；其中月日、时刻和地点来自原文，
年份来自通知的发布背景。乙同日问：“我是校外来宾，最晚什么时候报名？”
答复应保留严格条件：“须在2025年4月16日18:00前提交报名信息。”不能把它写成
“4月16日前”，也不能把18:00整当作已经满足“前”的要求。丙在4月17日问：
“我明天下午到可以吗？”此时“明天”是4月18日，但还需询问校内／校外身份；
若丙是校外来宾，必须核验是否已按时报名，不能因为活动尚未开始就承诺可以补报。

这些咨询使“理解通知”有了可检查的输出。词级分析可标作
“校外／来宾／须／报名”，其中“须”表达义务而不是已经完成报名；通知类别可为
“活动通知、时间变更”。实体记录包括“机器学习读书会”和“海棠楼201”，改期事件则
关联旧日期4月17日与新时间4月18日14:00，不能把两者当成两个独立活动。
下面给出几种文本与声音任务的输出规格，后文再把同一材料写成位置、字段和程序。

\begin{quote}
摘要：机器学习读书会改至2025年4月18日14:00在海棠楼201举行。校内师生可直接
参加，校外来宾须在4月16日18:00前报名。

译文：The Machine Learning Reading Group, originally scheduled for April 17,
has been moved to Friday, April 18, 2025, at 14:00 in Room 201, Haitang
Building. University students and staff may attend directly. External guests
must submit their registration information before 18:00 on April 16.
\end{quote}

译文中的年份使用已声明的背景；原文没有主讲人、费用和时长，译文也不补写。
朗读输出不能只交付上述文字：合成器应生成可播放的波形。例如把新时间读为
“二〇二五年四月十八日，星期五，下午两点”，地点读为“海棠楼二零一室”，截止条件读为
“四月十六日下午六点前”。这里给出的是教学用读出形式和验收规格，没有声称已生成
或测量真实音频。音频还须回听，检查数字、专名、停顿及是否漏读“前”。

人工逐篇阅读能够利用丰富背景，也能在信息不足时追问，但处理速度和一致性受人员、
工作量与操作规范影响。关键词规则执行快、结果容易追踪，例如看到“校外”和“报名”便
触发提醒；然而“校外来宾须报名”和“本次不接受校外报名”包含相似词语，规则若不处理
否定、作用范围和条件对象，就可能给出相反结论。由数据学习的模型可以覆盖更多表达，
仍会受到训练分布、输入截断和错误相关性的限制。自动预测不是免除任务定义，而是要求
把原本由人临场判断的对象、条件和验收方式写得更明确。

具体错误可以落在不同环节：关键词规则若取第一个日期，会把旧日期4月17日登记为新日期；
人工抄录若漏掉“18:00”，会改变校外来宾的报名条件；模型若把“我明天下午到”直接
映射成“欢迎参加”，则跳过了身份和报名状态。领域术语也会影响处理：
“机器学习”是活动主题，“读书会”是活动形式，不能仅因出现“学习”就归为课程报名。
这些错误分别要求时间关系、字段完整性、信息补充和类别规范，不能只靠扩大关键词表
或提高一个整体分数解决。

图~\ref{fig:nlp-one-text-many-tasks}把同一通知连接到若干输出。箭头表示原文可以被用于
哪些任务，并不表示分析、分类、抽取、生成和语音必须依次执行。一个端到端模型可以直接
从通知预测摘要；一个组合系统也可以先抽取事件记录，再根据记录回答问题。两条路线只要
交付对象相同，就应接受同一项最终约束，例如都不能把4月18日改写成4月17日。

\begin{figure*}[htbp]
  \centering
  \begin{tikzpicture}[
  every node/.append style={font=\figurefont\footnotesize},
  source/.style={book node,fill=FigureInput!12,draw=FigureInput,
    text width=38mm,minimum height=42mm,align=left},
  task/.style={book node,fill=FigureModel!10,draw=FigureModel,
    text width=23mm,minimum height=12mm,align=center},
  result/.style={book node,fill=FigureOutput!10,draw=FigureOutput,
    text width=65mm,minimum height=12mm,align=left}
]
  \node[source] (notice) at (21mm,0) {
    \textbf{教学通知}\\
    ML-0418\\
    2025年4月15日发布\\
    读书会改至4月18日\\
    14:00，海棠楼201。\\
    校外来宾须在\\
    4月16日18:00前报名。
  };
  \node[task] (analysis) at (67mm,54mm) {分析结构};
  \node[task] (classify) at (67mm,36mm) {类别};
  \node[task] (extract) at (67mm,18mm) {事实记录};
  \node[task] (summary) at (67mm,0) {摘要};
  \node[task] (translation) at (67mm,-18mm) {译文};
  \node[task] (answer) at (67mm,-36mm) {答案};
  \node[task] (speech) at (67mm,-54mm) {朗读};

  \node[result] (analysis-r) at (121mm,54mm)
    {“须报名”的义务对象是校外来宾};
  \node[result] (classify-r) at (121mm,36mm)
    {活动通知、时间变更};
  \node[result] (extract-r) at (121mm,18mm)
    {新时间：2025-04-18 14:00\\地点：海棠楼201};
  \node[result] (summary-r) at (121mm,0)
    {保留新时间、地点和精确报名条件};
  \node[result] (translation-r) at (121mm,-18mm)
    {Friday, April 18, at 14:00\\Room 201, Haitang Building};
  \node[result] (answer-r) at (121mm,-36mm)
    {“明天能来吗？”需日期、身份及报名状态};
  \node[result] (speech-r) at (121mm,-54mm)
    {“四月十八日，下午两点”\\输出为可播放声音，不能只交付文字};

  \draw[NNLine,line width=0.85pt] (notice.east) -- (48mm,0);
  \draw[NNLine,line width=0.85pt] (48mm,-54mm) -- (48mm,54mm);
  \foreach \name in {analysis,classify,extract,summary,translation,answer,speech} {
    \draw[book arrow] (48mm,0 |- \name.west) -- (\name.west);
    \draw[book arrow] (\name.east) -- (\name-r.west);
  }
\end{tikzpicture}

  \caption{同一份活动通知对应的不同语言任务。左侧为输入，中列为任务，右侧为输出。
  分析结果要符合标注规范，事实记录与摘要要对应原文，译文要保持意义，回答还可能依赖
  日期和业务状态，朗读须检查声音中的数字与专名。箭头表示材料用途，不表示任务执行顺序。}
  \label{fig:nlp-one-text-many-tasks}
\end{figure*}

具体应用还会增加通知之外的条件。咨询“我明天下午到可以吗”至少缺少提问日期和咨询者
身份；系统若已知当前是4月17日，才能把“明天”解释为4月18日。即便日期吻合，校外来宾
是否已经报名仍要查询业务状态。自然语言中的歧义、省略和多种合理表达，会把这些条件
带入任务定义。与其笼统要求系统“理解通知”，不如分别写清输入、输出、可用信息和
正确性依据。

\FloatBarrier
\section{自然语言处理的基本概念}
\label{sec:nlp-basic-concepts}

一项语言任务至少包含四个要素：处理的语言对象、要求产生的输出、允许使用的信息，以及
判断结果是否成立的依据。模型可以把输入变成内部向量，但向量本身通常不是应用要交付的
结果。用户需要的是标签、原文位置、事实记录、文本、程序或已经完成的行动；这些对象
必须能够还原、解析或检查。

\subsection{语言单位与上下文}
\label{sec:nlp-units-and-context}

语言材料可以按不同粒度观察。中文字符串“校外来宾须报名”按字符切为
“校／外／来／宾／须／报／名”，按本例词典切为“校外／来宾／须／报名”，
按一个教学子词词表则切为“校／外／来宾／须／报／名”。
\term{子词}（subword）是词表中可以拼接成较长词语的片段，不必是语言学上的词。
一个\term{词元}（token）是具体处理流程采用的离散输入单位，它可能对应一个字、一个词、
一个子词片段或标点。语言学意义上的“词”与模型词元因此不是同一概念。字节对编码
（byte pair encoding, BPE）、WordPiece和Unigram语言模型分别按合并、词表构造或概率
分割目标处理字符串，不保证上述三组结果相同。这些切分是教学构造，不冒充真实分词器的
输出。词元进入嵌入矩阵的接口见第~\ref{sec:transformer-input}节；本章不展开训练词表的算法。
字符组成句子，句子构成文档；对话还按说话者和发送时刻分成轮次。
句子、文档和轮次的边界决定哪些片段被一起输入，却不一定就是模型的特殊词元。

设原文为字符串$x$，使用从零开始、左闭右开的字符区间$[s,e)$表示一个
\term{跨度}（span），即从第$s$个位置取到第$e$个位置之前。本章按Unicode码点计数，
不是UTF--8字节数，也不是模型词元数。通知中“海棠楼201”的跨度为$[39,45)$，
长度为$45-39=6$；若单独输入“地点为海棠楼201”，局部跨度则为$[3,9)$。
分段记录必须同时保存局部位置和段落在完整文档中的起点。
左闭右开约定使相邻片段可以写成$[s,m)$与$[m,e)$，不会重复
包含边界字符。若系统先把全角字符转为半角、统一英文大小写或折叠连续空格，就必须保存
规范化位置到原文位置的映射；否则模型即使在规范化文本中找对片段，也可能无法在原通知
中高亮或回填。

下面用附加的混合书写片段演示映射，不把它当作通知的另一版本。原文为
“ＭＬ\ \ 海棠楼２０１\ \ ”，其中字母后、末尾各有两个ASCII空格；
依次作兼容规范化、英文小写化、连续空格折叠和首尾空格删除，得到
“ml 海棠楼201”。表~\ref{tab:nlp-units-offsets}的模型词元列忽略空格但保留其映射记录。

\begin{table}[htbp]
  \centering
  \begin{tabularx}{\linewidth}{@{}p{24mm}p{24mm}p{19mm}X@{}}
    \toprule
    原文及偏移 & 规范化及偏移 & 模型词元 & 变换 \\
    \midrule
    ＭＬ $[0,2)$ & ml $[0,2)$ & \code{ml} & 全角转半角，再转小写 \\
    两空格 $[2,4)$ & 一空格 $[2,3)$ & 无 & 多对一折叠 \\
    海棠楼 $[4,7)$ & 海棠楼 $[3,6)$ & 海棠楼 & 字符不变，偏移改变 \\
    ２０１ $[7,10)$ & 201 $[6,9)$ & \code{201} & 全角转半角 \\
    两空格 $[10,12)$ & 空串 $[9,9)$ & 无 & 删除，另记编辑记录 \\
    \bottomrule
  \end{tabularx}
  \caption{规范化前后的具体位置映射。每个区间左闭右开；词元来自教学词表，空格变换也
  必须记账，不能用一个全局偏移量替代整张映射。}
  \label{tab:nlp-units-offsets}
\end{table}

若模型选择规范化跨度$[3,9)$，合并第三、四行的原文区间后得到$[4,10)$，
从保留的原字符串切出“海棠楼２０１”。只保存规范化字符串无法恢复字母大小写、全角字形
和删除的空格；位置映射并不使有损变换可逆，原文与编辑记录必须保留。
一般规范化还可能把一个码点展开为多个码点，映射应允许多对多关系。

位置对应解决“结果落在原文哪里”，却没有解决“解释时可以看到什么”。
\term{上下文}（context）是当前预测允许参考的相关信息。它可以是同一句中的其他词语、
相邻句子、整篇文档、此前的对话轮次，也可以是检索到的外部资料。不同上下文提供的证据
并不等价。通知中的“现改至4月18日”依赖前面的原日期才能解释“改至”；咨询中的“明天”
依赖参考日期；“我已经报过名了吗”还依赖业务系统中的记录。

这些信息可以用四个小例子区分。“通知了老师的学生”可以指“老师的学生”，也可以指
“通知了老师”的学生，这是句内结构歧义；“小林联系了小周。她已报名”没有足够线索时，
“她”不能被武断归给最近的人名，这是跨句指代问题。若4月17日发来“时间改到明天下午”，
只能确定日期为4月18日，不能单凭“下午”推定14:00；旧通知和后续确认决定具体时刻。
对话“你是校内师生还是校外来宾？”“校外的”省略了主语和完整谓语，含义来自前一轮，
而非这三个字独自提供。

上下文范围必须在数据构造和部署时保持一致。训练样本若附带完整文档，部署时却只输入
一个被截断的句子，模型面对的是更少的信息。反过来，测试时若额外指定本应自行检索的
答案文档，或提供未来消息及人工测试答案，又会高估实际信息条件下的能力。下一章将具体讨论词法、
句法、语义、篇章与语用分析；本章只建立它们共同使用的单位、位置和上下文概念。

\subsection{任务输出与正确答案}
\label{sec:nlp-outputs-and-correctness}

语言任务的输出不只有文本。判断“咨询者是否需要报名”可以输出类别；估计咨询与通知的
相关程度可以输出数值；找出“海棠楼201”可以输出跨度；句法分析可以输出树，语义或
关系分析可以输出图；活动登记输出带字段的事实记录；摘要和译文输出新文本；数据库查询
输出程序；报名办理则以环境状态和回执作为结果。输出形式决定了哪些错误能够被直接检查。

仍以地点为例，跨度$[s,e)$命中原文，只说明位置正确。把该跨度写入
\code{venue}字段，还要检查字段模式、规范化规则和对应活动；摘要写成“活动在海棠楼举行”
虽然没有逐字复制原文，只要语义保持且没有引入无依据内容，仍可能合格。若系统声称“已为
你报名”，则不能用文字流畅或字段正确代替执行证据，必须检查报名记录是否创建成功。
位置命中、记录正确、语义等价、证据充分和行动完成是五种不同的验收对象。

把这些对象具体填入输出记录，可以得到如下约定。通知的多标签为
$\{\text{活动通知},\text{时间变更}\}$；咨询与通知的相关度可输出0到1之间的连续分数，
例如教学预测0.92，但这个分数并不自动是“答案正确概率”。地点输出为
\code{type=LOC, span=[39,45), text=海棠楼201}。依存树则使用约定的词切分：
“校内师生／可／直接／参加”，以“参加”为根，其余三个词的中心词均为“参加”，
依存标签依次为主语、助动词、状语。这是对通知片段的简化词级标注；
更换切分或标签规范以后，树的节点和标签也要一起改变。

活动记录不仅包含当前时间，还保留变更和出处。代码~\ref{code:nlp-event-record}中，
旧日期只记录到日，不能凭空补成某个旧时刻；新时刻的年份和时区来自已声明的背景。

\begin{codeblock}[label={code:nlp-event-record}]{活动与改期的教学记录}
event_id: ML-0418
name: 机器学习读书会
previous_date: 2025-04-17
starts_at: 2025-04-18T14:00:00+08:00
venue: 海棠楼201
registration:
  required_for: 校外来宾
  deadline_exclusive: 2025-04-16T18:00:00+08:00
source: notice-20250415
evidence:
  name: [0, 7]
  old_date: [10, 15]
  new_date: [21, 26]
  new_time: [30, 35]
  venue: [39, 45]
  deadline: [62, 72]
\end{codeblock}

程序输出还必须依赖已知数据库模式。补充设定表
\code{events(event\_id, starts\_at, venue)}，其主键为活动编号，时间字段使用带时区的
ISO格式字符串。问题“ML-0418什么时候、在哪里举行？”可以恢复为
代码~\ref{code:nlp-query}，在仅含上述活动的一行教学数据库上，返回
\code{2025-04-18T14:00:00+08:00}与“海棠楼201”。
它是只读查询，不会完成报名，也不会修改日历。

\begin{codeblock}[language=SQL,label={code:nlp-query}]{通知查询的程序输出}
SELECT starts_at, venue
FROM events
WHERE event_id = 'ML-0418';
\end{codeblock}

执行回执则来自工具而不是语言模型自行写出的文本。另设校外来宾甲已授权、资料完整，
工具在4月16日17:30完成报名，才允许返回如下记录：
\code{event\_id=ML-0418}、\code{registration\_id=R-001}、
\code{status=confirmed}、\code{committed\_at=2025-04-16T17:30:00+08:00}。
这是显式补充的业务状态，不是从通知推导的事实；验收还需用报名编号回查已持久化记录。
若工具只说“已接收请求”，或调用超时，系统只能报告待确认，不能照抄成功回执。

语言输出经常不止一个合理答案。若用途是面向所有潜在参加者发布简短提醒，则应保留
新时间、地点和报名条件。前节摘要符合要求；另一个合格摘要是
“请于2025年4月18日14:00到海棠楼201参加机器学习读书会，原4月17日安排取消。
校内师生无需报名；校外来宾请在4月16日18:00前提交报名信息。”
两者词序和措辞不同，但关键事实一致。若仅用于已报名人员的日历标题，可以另行允许
省略报名条件，不能把这种较窄用途的摘要当作通知摘要的通用标准。
多参考答案能够覆盖部分表达差异，却不可能枚举全部等价说法。因此，
生成任务通常需要结合表面匹配、语义比较、事实核验和人工判断，而不能把“与某条参考答案
逐字相同”当作唯一正确性定义。

标注本身也可能存在合理分歧。对“机器学习读书会”作实体识别时，标注规范可能把整个
短语记为活动名称，也可能只把“机器学习”标作主题。两种标注不能在同一数据集中任意混用，
否则评估会把规范差异误作模型错误。建立数据集时应固定边界、类型和冲突处理规则，并报告
标注者分歧；模型在测试集上复现标注规范，不等于它给出了脱离规范后唯一正确的语言分析。

有些问题在给定信息下没有答案。若只提供通知，询问“主讲人是谁”就应返回无法从材料确定，
而不是依据常见活动形式猜测。这里的“无答案”是相对于允许使用的材料而言；它既不表示
世界上不存在答案，也不表示模型永久不知道。任务规范需要区分无答案、未知类别、输入
无效和工具失败，因为它们对应不同的后续处理。

\subsection{数据、监督与信息条件}
\label{sec:nlp-data-supervision-conditions}

原始文本只给出语言材料，训练任务还需要把材料与目标联系起来。分类样本可以写成
$(x,y)$，其中$x$是通知或咨询，$y$是类别；序列标注样本为每个处理位置提供BIO等标签；
关系与事件样本记录提及位置、类型、论元及证据；翻译样本由源文与目标文组成；问答样本
则可能同时包含问题、证据、答案和可回答性标记。行动任务还要记录允许的工具、动作参数、
环境反馈与完成状态。格式不同的样本提供不同监督，不能只按“都有文本”把它们互换。

代码~\ref{code:nlp-supervision-records}列出可实际存储的样本。这里的
\code{doc}指向上述80字符原文，\code{LOC}是地点类型，BIO中的
\code{B-LOC}表示地点的第一个单位，\code{I-LOC}表示其后续单位，
\code{O}表示当前不属于该类实体。关系样本使用命名实体间的“举办于”关系，
事件样本则有“旧日期”“新日期”等不同角色，不能仅把字段排成无类型的字符串。

\begin{codeblock}[label={code:nlp-supervision-records}]{同一材料的监督样本}
分类:
  input: {doc: notice-20250415}
  labels: [活动通知, 时间变更]
序列标注:
  text: 地点为海棠楼201
  units: [地, 点, 为, 海, 棠, 楼, 2, 0, 1]
  tags: [O, O, O, B-LOC, I-LOC, I-LOC,
         I-LOC, I-LOC, I-LOC]
  document_offset: 36
关系:
  head: {span: [0, 7], type: EVENT_NAME}
  tail: {span: [39, 45], type: LOC}
  relation: held_at
  evidence: [0, 46]
事件:
  type: schedule_change
  trigger: [19, 21]
  old_date: [10, 15]
  new_date: [21, 26]
  new_time: [30, 35]
翻译:
  source: 地点为海棠楼201。
  target: The venue is Room 201, Haitang Building.
问答:
  question: 活动在哪里举行？
  evidence: {doc: notice-20250415, span: [39, 45]}
  answer: 海棠楼201
  answerable: true
\end{codeblock}

行动样本需要展开时间和反馈。以刚才的成功报名为例，轨迹为：
用户在截止前授权提交资料；助手调用\code{register(ML-0418, guest-A)}；
工具返回\code{confirmed, R-001}；助手调用\code{get\_registration(R-001)}；
工具返回该活动与该来宾的已确认记录；助手才告知“已报名，编号R-001”。
训练记录应分别保存\code{role}、\code{content}、\code{tool\_call}、
\code{observation}和最终成功标记。若只把最后一句当作参考回复，模型学到的是
如何表述成功，缺少何时可以宣告成功的监督。

监督信息可以由人工标注、词典与规则、已有模型或数据自身的关系产生。规则匹配可以快速
构造候选标签，但规则错误会进入训练数据；\term{远程监督}（distant supervision）借助
外部知识记录与文本对齐产生标签，其标签只在对齐假设成立时可靠；\term{伪标签}
（pseudo-label）是模型暂时赋给未标注样本的预测，不会因为参与训练就变成真实标注。
指令示例则提供“输入该怎样解释、输出该遵循什么格式”的实例；它可能作为微调样本，
也可能只放在当前提示中，不能把二者都称为发生了参数学习。自监督目标可从原文遮蔽
“海棠楼201”再要求恢复，提供词元预测目标，却不直接给出“举办地点”的业务字段。
这些信息来源的比较见第~\ref{chap:machine-learning-paradigms}章，
尤其第~\ref{sec:self-supervised-learning}、\ref{sec:weakly-supervised-learning}
和~\ref{sec:transfer-learning}节。

例如人工人员按统一规范标出$[39,45)$，能同时提供边界和类型；
地点词典命中“海棠楼”只提供$[39,42)$，未覆盖房间号。若外部表记录活动在海棠楼201，
远程监督把提到该活动和地点的文档标为正例，仍可能误标讨论旧安排的句子；
现有模型产生的伪标签又可能继承它对新旧日期的混淆。监督的来源、置信度与是否经人工
复核应当进入记录，不能只留下一个看似确定的标签。

评价样本必须与模型选择过程分开。若训练、提示编写、阈值选择或错误分析已经读取测试
答案，测试结果便不再代表未见数据。语言数据尤其容易出现隐蔽重复：同一活动通知可能以
网页、邮件和聊天转发三种形式出现，随机按句子切分会让近乎相同的事实同时进入训练集与
测试集。根据部署问题，可以按文档、来源、实体或时间切分。例如要预测未来通知，按时间
留出比随机切分更接近使用条件；要检验新活动和新机构，则应避免同一实体跨集合泄漏。

可把这几种切分落实成可重复的协议：按文档切分时，原通知及删改标点的转发版本归为一组；
按来源切分时，以发布机构或站点为组留出，而非把同一机构的网页和邮件视为独立来源；
按实体切分时，\code{ML-0418}的所有通知、咨询和记录必须在同一集合；
按时间切分时，以4月1日至14日训练、15日验证、16日至18日测试作为教学边界，
并只给每个测试请求当时已发布的资料。实际数据须有足够多活动，本章单个活动只能
示范协议，不能独自构成有效测试集。
同模板新通知还可能复制机构固定措辞，故须另报模板相同和模板未见两组；
按日期切分也不自动消除这种相似。开放资料问答可以访问测试问题所需的原始通知，
但不能把人工测试答案、评分解释或未来更新混入检索库。

除了切分方式，还要记录语言、领域、资源规模、文档长度、资料时点和允许访问的信息。
“低资源”表示目标语言或任务可用的标注、平行文本或工具有限，并不是某个模型的固有属性；
“未见类别”也只有在训练类别与测试类别的关系被明确后才成立。若问答系统在2025年的
资料库上回答2026年的活动，错误可能来自资料时点，而不是阅读器没有识别句子。任务条件
应让这类来源能够被区分。

\section{自然语言处理的研究内容}
\label{sec:nlp-research-content}

前一节从任务定义说明“要处理什么”。自然语言处理研究还需要回答三个彼此制约的问题：
怎样组织任务之间的关系，已有模型怎样把输入变成所需输出，以及用什么证据判断系统适合
实际使用。任务规定目标，实现给出计算路径，评价则检查路径产生的结果；三者不能相互
替代。

\subsection{文本分析、文本处理与扩展任务}
\label{sec:nlp-analysis-processing-extended}

文本分析描述语言材料中的单位、结构、意义和上下文关系。词法分析确定局部单位与属性，
句法分析表示成分或依存关系，语义分析刻画词义、论元和形式含义，篇章分析连接跨句实体
与关系，语用分析则考察说话目的、共同背景与交流行为。分析结果可以直接交付，也可以
成为其他任务的输入。例如关系抽取可以利用依存路径缩小候选关系范围，但也可以由编码器
直接为实体对评分。显式分析提供可检查的中间结构，同时会传播切分或解析错误；直接预测
减少中间接口，却仍需用数据、解码约束或后验检查保证输出结构合法。

文本处理任务按使用目的选择所需信息。分类为整段文字或文本对赋予类别，抽取把原文事实
整理成实体、关系与事件，组织任务按相似性、主题或相关性排列材料，生成与转换产生摘要、
译文、改写或数据描述，问答和对话围绕用户需求选择证据并形成回复。这些任务并不按固定
流水线排列。垃圾信息分类未必需要显式句法树，文档级事件记录却往往要处理跨句指代。

一个完整应用可以组合多个任务。活动登记系统先从通知抽取活动、时间、地点和条件，再按
数据库模式规范化字段，最终以写入记录是否通过约束检查验收。资料咨询系统先检索相关
通知，再读取证据并回答问题，检索命中与答案受支持需要分别检查。改期办理还要维护对话
状态、识别用户授权、调用日历或报名工具，并用环境回执确认操作完成。组合系统中的中间
接口方便定位错误，但最终成功不能由某个模块的高分代替。

对通知事实登记，输入是原文及发布时刻，中间产物是带跨度的实体、改期事件和
代码~\ref{code:nlp-event-record}中的规范化字段。记录合并以活动编号关联已有记录，
将旧日期保留为历史，新日期成为当前安排；输出是数据库中的一条当前记录及其来源。
最终检查活动没有被重复创建、字段与原文对应、旧日期未覆盖新日期。
对资料咨询“原来的日期还有效吗”，检索输出通知编号及相关片段，阅读输出
“原定4月17日，现改至4月18日”，回答核验再检查原文是否明确表达了取代关系。
最终答复是“原4月17日安排已改为4月18日”，不是仅凭检索命中就说“是”。

改期办理则补充一项明确的业务设定：用户自己的日历仍登记4月17日，且用户已授权同步
到新时间。对话状态应含活动编号、目标时间、目标日历和授权标记；
工具参数为\code{update\_event(ML-0418, 2025-04-18T14:00:00+08:00)}。
观察若返回\code{updated, revision=2}，系统还应读取该版本并核对日期和地点，
再输出“日历已更新，版本2”。若用户说“只查一下，别改”，授权标记应为假，
流程停在回答阶段。通知抽取正确与日历确实改期是两项不同的验收结果。

这也解释了本部分的章节安排：下一章交付单位、结构和意义分析；信息抽取章交付可追溯
事实；分类、匹配与组织章处理单篇、文本对和集合；生成与转换章交付新文本；
问答与对话章增加证据和交流历史；求解与执行章增加程序、动作及完成证据；
语音章再增加声音和时间条件。这是按交付对象分工，不是要求所有应用都逐章运行。

自然语言还可以连接可执行对象与声音。语义解析把问题转换为数据库查询或其他形式表达；
程序生成和语言智能体把要求变成代码或环境动作；语音识别、合成和口语对话在文字之外
增加声学输入输出与实时性条件。本部分讨论这些任务中由语言决定的接口和评价。页面版面、
图像、视频及其与文字的联合证据属于后续图像处理部分。

\subsection{已有模型怎样形成任务输出}
\label{sec:nlp-model-output-interfaces}

不同模型可以共享一个任务，同一模型也能通过不同输出接口服务多个任务。为了比较实现，
需要追踪模型为哪个候选对象计算得分，怎样从得分恢复结果，以及恢复后还要检查什么。
表~\ref{tab:nlp-model-output-interfaces}按这条路径概括几类常见接口。

\begin{table*}[htbp]
  \centering
  \begin{tabularx}{\linewidth}{@{}p{31mm}p{35mm}p{42mm}X@{}}
    \toprule
    实现接口 & 得分或预测对象 & 输出恢复 & 仍需检查的约束 \\
    \midrule
    TF--IDF加朴素贝叶斯或线性分类器
      & 文档与每个候选类别
      & 选择最大得分类别，或按验证集阈值输出多标签
      & 特征词表、未知词、类别代价与概率校准 \\
    编码器加分类头或跨度头
      & 整段类别，或每个起止位置对
      & 类别取最大值；跨度按边界得分和合法性筛选
      & 截断、重叠跨度及词元到原文位置的映射 \\
    HMM或线性链CRF
      & 位置标签与相邻标签转移
      & 用Viterbi算法恢复整体得分最高的合法序列
      & 标签体系、空序列、起止状态及平局规则 \\
    编码器--解码器或自回归语言模型
      & 给定输入与已有前缀后的下一词元
      & 贪心、束搜索或采样，遇终止符或长度上限停止
      & 停止条件、重复、格式解析、事实与任务约束 \\
    检索器、阅读器与工具接口
      & 资料候选、证据片段、答案或动作参数
      & 组合检索结果，形成回答或提交结构化调用
      & 来源时点、权限、参数模式与环境回执 \\
    \bottomrule
  \end{tabularx}
  \caption{常见模型接口怎样形成语言任务输出。表中只比较输入到输出的任务接口；模型训练、
  概率推断与优化细节见前卷对应章节。}
  \label{tab:nlp-model-output-interfaces}
\end{table*}

以“是否含报名规定”为二分类任务，TF--IDF把词语频率及其跨文档稀有程度变成特征。
朴素贝叶斯用类别先验和词项似然形成类别得分，线性分类器用
$s_k=\bm{w}_k^\top\bm{x}+b_k$为类别$k$评分，$\bm{x}$为文档特征。
例如两类得分为$(1.8,-0.2)$时取前一类“含报名规定”，却不能由此断定所有人都必须报名。
词表和逆文档频率只在训练语料上拟合，分类阈值在验证集上选定。
这里的关键输出是类别，不是TF--IDF向量。
分类模型与评价见第~\ref{chap:classification}章。若改用BERT一类编码器，可取整段表示接
分类头，也可为每个起止位置对计算实体跨度得分；Transformer的信息流和预训练接口见第
\ref{chap:transformer}章。

跨度头若对规范化词元“海棠楼”“201”分别给出最高的起、止分数，输出的是连续词元区间；
需要经过表~\ref{tab:nlp-units-offsets}的映射才能恢复原文字符跨度。
不能把模型里的第3个词元误作原文第3个字符，也不能独立选出终点在起点之前的两个高分位置。

序列标注要求相邻位置的答案共同合法。例如BIO标签中的\code{I-LOC}不应无条件出现在
地点片段开头。HMM或线性链CRF为局部标签及其转移评分，再由Viterbi算法恢复整体最优
序列；逐位置分别取最大分数可能拼出不合法或整体得分更低的路径。HMM推断见第
\ref{sec:pgm-dynamic-discrete}节，CRF定义与解码见第
\ref{sec:pgm-energy-models-crf-definition}节。

例如“海／棠”两个位置的局部最高标签为\code{I-LOC, I-LOC}，逐点取最大值会以
\code{I-LOC}开头。带起始约束的Viterbi应在所有合法序列中比较发射分和转移分，
可能恢复\code{B-LOC, I-LOC}；被选择的不是各位置独立最优标签的拼接。

条件生成把目标文本写成词元序列$y_{1:T}$，给定输入$x$后按链式法则建模
\begin{equation}
  p(y_{1:T}\mid x)
  =
  \prod_{t=1}^{T}p(y_t\mid y_{<t},x).
  \label{eq:nlp-conditional-generation}
\end{equation}
训练可以增加参考输出在这些条件分布下的概率，使用时则要选择具体词元并决定何时停止。
束搜索保留若干高分前缀，采样按分布产生不同候选；二者改变搜索与输出分布，不会自动保证
日期正确或摘要完整。分类和抽取也可以改写成生成标签或结构化字符串，但生成后仍要解析
合法类别、恢复原文证据，并按原任务标准验收。

例如输入“只返回通知类别”，生成器输出\code{时间变更}后，还要确认它属于许可标签集合；
输入“抽取地点并给出证据”，输出\code{venue=海棠楼201; span=[39,45)}后，
解析器应检查字段名和类型，再验证\code{venue}与原文切片一致。
字符串能够解析，只证明格式有效；若输出“海棠楼301”，格式检查仍可能通过，
原文核验却必须失败。工具参数同理：即使时刻格式合法，报名时间晚于截止点也不能提交。

外部检索和工具调用进一步改变可用信息。检索器把问题与资料候选匹配，阅读器或生成器
使用返回片段形成答案；工具接口则把语言输出解析为函数名与参数，在权限检查后交给环境。
检索到一篇相关通知不等于回答已经受到支持，生成一个外观合法的调用也不等于操作成功。
来源、参数和回执分别提供不同阶段的证据。

\subsection{评价依据与应用比较}
\label{sec:nlp-evaluation-and-comparison}

评价从输出对象出发。类别可以比较预测标签与参考标签，跨度可以检查边界和类型，树或图
需要比较结构，事实记录需要核对字段与来源，新文本需要检查意义、事实与表达，程序和行动
则可以运行测试或查看环境状态。一个易于计算的替代指标只有在与任务要求建立了明确关系
时才有解释力。

以判断咨询者“需要报名”为正类的二分类为例，构造50个已经明确身份的测试样本，
按“真实类别为行、预测类别为列”得到混淆矩阵
\[
 \begin{array}{c|rr}
  & \text{预测需报名} & \text{预测无需报名}\\ \hline
  \text{真实需报名} & 18 & 6\\
  \text{真实无需报名} & 2 & 24
 \end{array}.
\]
真正例、假正例、假负例和真负例分别为
$\mathrm{TP}=18$、$\mathrm{FP}=2$、$\mathrm{FN}=6$和$\mathrm{TN}=24$。准确率、精确率、
召回率与F1分别为
\begin{align}
  \mathrm{Accuracy}
  &= \frac{\mathrm{TP}+\mathrm{TN}}
           {\mathrm{TP}+\mathrm{FP}+\mathrm{FN}+\mathrm{TN}}
   = \frac{42}{50}=0.84, \\
  \mathrm{Precision}
  &= \frac{\mathrm{TP}}{\mathrm{TP}+\mathrm{FP}}
   = \frac{18}{20}=0.90, \\
  \mathrm{Recall}
  &= \frac{\mathrm{TP}}{\mathrm{TP}+\mathrm{FN}}
   = \frac{18}{24}=0.75, \\
  F_1
  &= \frac{2\,\mathrm{Precision}\,\mathrm{Recall}}
           {\mathrm{Precision}+\mathrm{Recall}}
   \approx 0.818.
  \label{eq:nlp-classification-metrics}
\end{align}
这组数值说明模型给出的正例较少出错，却漏掉了四分之一的真实正例。若任务是提醒校外
来宾报名，漏报和误报的代价可能不同，不能只按准确率选模型。多类别任务中的宏平均先对
每个类别分别计算指标再等权平均，能显示小类别表现；微平均先汇总各类别计数，更受大类别
影响。两种平均回答的问题不同，应与逐类别结果一同报告。

在这张矩阵中，把负类当作当前关注类别后，
$P_-=24/(24+6)=0.8$，$R_-=24/(24+2)=12/13$，
$F_{1,-}=48/56=6/7$；正类$F_{1,+}=36/44=9/11$。
两类等权平均得到
\begin{equation}
 F_{1,\mathrm{macro}}
 =\frac12\left(\frac9{11}+\frac67\right)
 =\frac{129}{154}\approx0.837662.
 \label{eq:nlp-macro-f1}
\end{equation}
微平均则把两个类别分别按一对其余计数后汇总：
$\sum_k\mathrm{TP}_k=42$，$\sum_k\mathrm{FP}_k=\sum_k\mathrm{FN}_k=8$，
故$P_{\mathrm{micro}}=R_{\mathrm{micro}}=F_{1,\mathrm{micro}}=42/50=0.84$。
单标签分类且纳入全部类别时微F1等于准确率；多标签或只取部分类别时不能套用这个等式。
宏F1是各类F1的平均，不是先平均精确率、召回率再求一次调和平均。

实体抽取的严格匹配通常要求起止位置和类型全部一致。“海棠楼”与参考跨度“海棠楼201”
发生部分重叠，但不能直接写入完整地点字段；部分匹配可以用于错误分析，不应在未说明规则
时替代严格结果。设唯一预测为\code{LOC:[39,42)}，唯一参考为\code{LOC:[39,45)}，
严格实体计数是$\mathrm{TP}=0,\mathrm{FP}=1,\mathrm{FN}=1$，F1为0。
若另报同类型字符重合诊断，交集长度为3，预测长度3，参考长度6，得到精确率1、
召回率$1/2$、F1为$2/3$。这是明确约定的字符级部分匹配，不是严格实体F1；
多个预测竞争同一参考时，还须先规定一对一匹配，防止重复计分。

句对相似度输出连续分数时，可以比较分值与排序。表~\ref{tab:nlp-correlation}固定第一句
为“活动改到4月18日”，参考分数$a_i$和模型分数$b_i$均取0到5，数值只是教学构造。

\begin{table}[htbp]
 \centering
 \begin{tabularx}{\linewidth}{@{}Xrrrr@{}}
  \toprule
  另一句 & $a_i$ & $b_i$ & $a_i$的秩 & $b_i$的秩\\
  \midrule
  地点是海棠楼201 & 1 & 1 & 1 & 1\\
  活动时间有调整 & 2 & 3 & 2 & 3\\
  活动安排在4月18日 & 3 & 2 & 3 & 2\\
  活动日期改为4月18日 & 4 & 5 & 4 & 4\\
  \bottomrule
 \end{tabularx}
 \caption{句对连续评分与排序的构造例子。秩从低分到高分排列，无并列分数。}
 \label{tab:nlp-correlation}
\end{table}

\term{Pearson相关系数}衡量两组分值的线性一致程度。这里$\bar a=2.5$、
$\bar b=2.75$，中心化后的交叉乘积之和为5.5，两组平方和为5和8.75，故
\begin{equation}
 r=\frac{\sum_i(a_i-\bar a)(b_i-\bar b)}
 {\sqrt{\sum_i(a_i-\bar a)^2\sum_i(b_i-\bar b)^2}}
 =\frac{5.5}{\sqrt{5\times8.75}}\approx0.831522.
 \label{eq:nlp-pearson}
\end{equation}
\term{Spearman秩相关系数}则对秩计算Pearson相关。没有并列秩时，令$d_i$为两组秩差，
可写为$\rho=1-6\sum_i d_i^2/[n(n^2-1)]$；本例$d_i=(0,-1,1,0)$，
$n=4$，所以$\rho=1-12/60=0.8$。有并列分数时应赋平均秩后计算相关，
不能不加修正地使用上述无并列简式。Pearson对整体平移和正比例缩放不敏感，
Spearman还忽略保持排序的非线性变换，二者都不能证明分值已校准或某个分类阈值合适；
任意一组分数为常量时相关系数无定义。

结构、生成、语音和程序任务各有专门指标。\term{无标签依存准确率}
（unlabeled attachment score, UAS）检查中心词是否正确，
\term{带标签依存准确率}（labeled attachment score, LAS）要求中心词与关系标签同时正确。
是否排除标点是两者都需要另行说明的评价选项，与UAS／LAS的区别无关。
Smatch比较抽象意义表示（AMR）图中的三元组对应，BLEU、ROUGE与SARI分别
从不同角度比较翻译、摘要或简化结果，词错误率（WER）与字符错误率（CER）检查转写编辑
距离，\code{pass@k}估计多个程序候选中至少一个通过测试的概率，任务成功率则检查环境
目标是否达到。这些指标的分母和遗漏内容各不相同，不能相加成一个没有定义的“综合分”。

\begin{table*}[htbp]
  \centering
  \begin{tabularx}{\linewidth}{@{}p{19mm}p{28mm}p{35mm}p{35mm}X@{}}
    \toprule
    输出对象 & 通知任务例子 & 直接正确性依据 & 常用替代指标 & 指标可能遗漏的要求 \\
    \midrule
    类别 & 是否需要报名 & 任务规范下的参考类别 & 准确率、宏／微F1 & 类别代价、阈值与校准 \\
    跨度 & 地点“海棠楼201” & 原文边界与实体类型 & 严格或部分匹配F1 & 规范化值和实体身份 \\
    结构 & 通知的依存或事件图 & 合法结构与参考关系 & UAS、LAS、Smatch & 标注体系差异及下游用途 \\
    事实记录 & 活动时间、地点与条件 & 字段值、模式和原文证据 & 字段级精确率与召回率 & 字段间一致性和资料时点 \\
    新文本 & 摘要、译文或答复 & 意义保持、覆盖、事实与用途 & BLEU、ROUGE、语义相似度 & 无参考表达、事实错误与风格约束 \\
    程序 & 查询活动数据库 & 语法、执行结果与目标语义 & 精确匹配、执行准确率 & 偶然同结果和副作用 \\
    环境行动 & 提交报名 & 授权后的目标状态与回执 & 任务成功率 & 用户负担、权限和不可逆后果 \\
    \bottomrule
  \end{tabularx}
  \caption{输出对象与评价证据。替代指标只覆盖直接正确性的一部分；完整应用还要固定输入
  条件、资源预算和失败处理方式。}
  \label{tab:nlp-output-evidence}
\end{table*}

自动指标适合稳定、批量地比较明确对象，人工评价可以判断用途、自然程度和未被形式规则
覆盖的错误，模型评审则能降低部分开放文本评价成本。三者都不是无条件真值。一条摘要
“活动改到4月17日下午”可能语言流畅，与参考文本也有大量词语重合，却改错了关键日期。
人工评价会受到说明、专业程度和上下文影响；模型评审还可能偏好靠前、较长或与自身表达
相似的答案。使用模型评审时，应交换候选顺序、固定评价量表，并以可独立核查的事实或人工
抽样校验其判断；还应固定评审模型、版本、提示和可见资料。MT-Bench与Chatbot Arena
的评价分析讨论了位置、冗长及自我偏好，不能把偏好胜出解释为事实正确
（\citeyear{nlp-intro-zheng2023judge}；完整来源见本章文献说明）。

固定测试集的平均分也可能掩盖系统性错误。CheckList把语言能力拆成最小功能测试、不变性
测试和方向性期望，用否定、实体替换、数量或语序变化检查模型是否依赖表面捷径
\scite{ribeiro2020checklist}
Dynabench进一步让人与模型迭代构造困难实例，观察静态测试
之外的新失败（\citeyear{kiela2021dynabench}）。这些方法能发现反例，却不能证明未测试的语言
现象都已覆盖。

可以围绕通知设计这样的行为检查：把“须报名”改为“不用报名”，相关类别必须改变；
把房间201一致地替换为301，返回地点必须同步改变；只交换“地点”和“时间”两分句的次序，
当前活动记录应不变；把报名人数从2改为3时，涉及数量的答案应相应变化。
“校外身份＋明天参加＋尚未报名”是需要单独保留的组合，不能由三个单项测试代替。
若两份通知对新日期冲突，应展示来源和版本，不能默选一个日期。把部署中发现的新失败
加入下一轮困难实例时，应单独保留未参与修订的测试集，不能把反复修过的同一批题当作
独立泛化证据。

跨任务和跨语言比较还要保留条件。HELM（\citeyear{nlp-intro-liang2023helm}）
把场景、适配方式和多项指标放在统一评价流程中，保存提示与输出以便复查；
XTREME（\citeyear{hu2020xtreme}）把分类、结构标注和检索等任务放入多语言评价，
逐任务与逐语言结果不能被单个总分抹平。Belebele（\citeyear{bandarkar2024belebele}）
使用平行篇章比较多种语言变体的阅读理解，Global MMLU（\citeyear{singh2025globalmmlu}）
则指出翻译题目还可能携带源语言与文化背景假设。因此，语言覆盖不仅是把同一
句子翻译成更多语言，还要检查目标用户、自然表达、知识背景与评价资料是否匹配。

完整应用的比较应同时记录质量、延迟、标注和调用成本、用户负担及维护要求。错误分析可
沿输入处理、检索、预测、解码、核验和执行逐段定位，但局部模块分数与最终任务成功需要
分别报告。阈值和提示在验证集上选择，独立测试集只用于最终评价；上线后若语言、领域、
资料时点或用户行为改变，还要重新检查原评价是否仍有代表性。

例如在验证集上比较报名提醒阈值0.4与0.6，按事先约定的漏报成本选定其一，然后冻结阈值，
在独立测试集上报告混淆矩阵。比较两个完整系统时，应统一是否允许查报名库、
资料截止时刻和最多工具调用次数，并分别报告校内／校外、已报名／未报名、
明确日期／相对日期的结果。一个系统多追问一轮后成功，另一个无追问但误答，
应同时记录成功与交互成本，不能只把成功率或平均时延单独当作全部结论。

\section{自然语言处理的研究历史}
\label{sec:nlp-research-history}

自然语言处理并不是沿一条模型时间线前进。语言结构、文本判定、事实记录、文本转换、
证据问答、程序行动和语音交互长期并行；近年的变化常表现为输入范围扩大、输出对象变复杂，
或评价条件更接近真实使用。图~\ref{fig:nlp-task-history}按本书后续任务组织若干代表节点，
连线表示问题上的延续关系，不表示后一个工作取代前一个工作。

本节采用文献地图所选版本的年份：已列出会议或期刊版本的工作按正式发表年，
只引用预印本的工作明确标记为预印本。例如果干模型早已公开，Whisper仍按ICML 2023年
论文定位，Tacotron 2按ICASSP 2018年论文定位；它们分别在2022年、2017年已有预印本。
这不表示任务或思想始于该年。完整出处见本章末的文献说明。

\begin{figure*}[htbp]
  \centering
  \begin{tikzpicture}[
  x=13.4mm,y=1mm,
  every node/.append style={font=\figurefont\footnotesize},
  event/.style={fill=white,inner sep=1.5pt,align=center,text=BookInk},
  guide/.style={draw=BookRule,densely dotted,line width=0.5pt}
]
  \node[anchor=west,align=left,font=\figurefont\footnotesize,text=BookMuted]
    at (-1.85,16) {2017年前的背景：规则与词典、HMM／CRF标注、词袋分类、pLSA／LDA};
  \foreach \year [evaluate=\year as \x using int(\year-2017)] in
    {2017,2018,2019,2020,2021,2022,2023,2024,2025,2026} {
    \draw[guide] (\x,3) -- (\x,-125);
    \node[anchor=south,text=BookMuted] at (\x,5) {\year};
  }
  \foreach \name/\y in {
    {结构与事实}/-8,
    {文本判定}/-30,
    {文本转换}/-52,
    {证据与交流}/-74,
    {程序与行动}/-96,
    {语音}/-118
  } {
    \node[anchor=east,font=\figurefont\footnotesize\bfseries]
      at (-0.42,\y) {\name};
    \draw[BookRule,line width=0.8pt] (0,\y) -- (9,\y);
  }

  \foreach \x/\y in {0/-8,2/-8,3/-8,5/-8,8/-8} {
    \fill[FigureInput] (\x,\y) circle (1.4pt);
  }
  \node[event,above] at (0,-7) {共指};
  \node[event,below] at (2,-9) {DocRED};
  \node[event,above] at (3,-7) {UD v2};
  \node[event,below] at (5,-9) {MAVEN-ERE};
  \node[event,above] at (8,-7) {EventRelBench};

  \foreach \x in {0,1,2,3} {\fill[FigureModel] (\x,-30) circle (1.4pt);}
  \node[event,above] at (0,-29) {STS};
  \node[event,below] at (1,-31) {MultiNLI};
  \node[event,above] at (2,-29) {PAWS};
  \node[event,below] at (3,-31) {ANLI／CLUE};

  \foreach \x in {0,1,2,5,7,9} {\fill[FigureOutput] (\x,-52) circle (1.4pt);}
  \node[event,above] at (0,-51) {JFLEG};
  \node[event,below] at (1,-53) {XSum};
  \node[event,above] at (2,-51) {STACL};
  \node[event,below] at (5,-53) {NLLB（预）};
  \node[event,above] at (7,-51) {WMT24};
  \node[event,below] at (9,-53) {DiscoX};

  \foreach \x in {1,3,8} {\fill[BookAmber] (\x,-74) circle (1.4pt);}
  \node[event,above] at (1,-73) {SQuAD 2.0};
  \node[event,below] at (1,-75) {HotpotQA};
  \node[event,above] at (3,-73) {RAG};
  \node[event,below] at (8,-75) {$\tau^2$-Bench（预）};

  \foreach \x in {4,6,7} {\fill[BookTheorem] (\x,-96) circle (1.4pt);}
  \node[event,above] at (4,-95) {GSM8K（预）};
  \node[event,below] at (4,-97) {HumanEval（预）};
  \node[event,above] at (6,-95) {ReAct};
  \node[event,above] at (7,-95) {SWE-bench};
  \node[event,below] at (7,-97) {WebArena};

  \foreach \x in {1,4,6,7,9} {\fill[BookAmber] (\x,-118) circle (1.4pt);}
  \node[event,above] at (1,-117) {Tacotron 2};
  \node[event,below] at (4,-119) {CoVoST 2};
  \node[event,above] at (6,-117) {Whisper};
  \node[event,below] at (6,-119) {SeamlessM4T（预）};
  \node[event,above] at (7,-117) {Moshi（预）};
  \node[event,below] at (9,-119) {$\tau$-Voice（预）};
\end{tikzpicture}

  \caption{2017至2026年若干语言任务问题的扩展。节点按本节声明的文献版本定位，
  “预”表示预印本；相邻标签错开摆放，定位点仍对应年份。HotpotQA与SQuAD 2.0均在2018年，
  Whisper按正式论文记为2023年。图中仅抽取代表节点，完整线索见正文；连线表示本文讨论的
  问题联系，不表示替代关系或未经核验的直接继承。}
  \label{fig:nlp-task-history}
\end{figure*}

\FloatBarrier
\subsection{语言结构与文本判定的研究}
\label{sec:nlp-history-analysis-judgment}

在近十年之前，规则与词典已经给出了可检查的边界和类别；HMM把标签序列与观测联系起来，
CRF直接对给定文本下的标签序列建模。词袋分类按文档中的词项统计判定类别，概率主题模型
则在无类别标注的文档集合中估计共同主题。检索和词袋方法的背景可参见
\citetitle{manning2008ir}，结构化标注可参见\citetitle{pgm-energy-models-sutton2012}。
这些接口持续与神经表示共存：词典仍决定候选覆盖，结构解码仍检查标签合法性，
上下文表示不会自动替代输出规范。

句法的前后对照是从显式语法规则及其概率，转向由上下文编码器给候选成分跨度评分，
但两者都要恢复合法树。Kitaev与Klein的成分解析工作
（\citeyear{nlp-intro-kitaev2018}）体现了编码器与结构恢复的分工。
共指的对照则是从先识别提及再选择先行词，转向同时为候选跨度和先行词关系评分；
端到端共指（\citeyear{lee2017coreference}）把提及识别与聚类联系起来，
并没有消除跨句省略或身份混淆。语义角色研究进一步比较“谓词已给定”与“谓词、论元共同
预测”的监督条件（\citeyear{nlp-intro-he2018}）。
UD v2（\citeyear{nivre2020universal}）用共享规范组织多语言的单位、词性和依存关系；
共享标注便于比较，不表示各语言的词序、形态或省略规律相同。

分类的一个对照是“通知含哪些词”与“词在什么关系中出现”。词袋判定容易把
“甲通知乙”和“乙通知甲”看成相似输入；上下文编码可以利用顺序，但是否真的用了顺序
仍须通过任务验证。STS（\citeyear{nlp-intro-cer2017}）要求连续相似度评分；
MultiNLI（\citeyear{nlp-intro-williams2018}，预印本2017年）在前提与假设间区分
蕴含、矛盾与中性，并增加体裁匹配／不匹配测试；
PAWS（\citeyear{nlp-intro-zhang2019paws}）以高词汇重叠而角色或语序不同的句对检验
词袋捷径。ANLI（\citeyear{nlp-intro-nie2020}）增加人与模型迭代构造的困难实例，
CLUE（\citeyear{nlp-intro-xu2020clue}）增加中文多任务条件，不能用它们的一个综合分
代替逐项语言能力判断。

主题发现的对照是另一种监督变化。pLSA（\citeyear{pgm-static-directed-hofmann1999}）
和LDA（\citeyear{blei2003}）用文档词项共现估计主题与文档混合关系；
上下文主题模型（\citeyear{nlp-intro-bianchi2021}）再把预训练文档表示作为推断信息。
例如含“报名、地点、改期”的通知可与另一篇活动安排共同形成主题，但主题编号不是人工
定义的“活动通知”类别，也不自动给出“应当如何回复”。主题连贯性和多样性仍需评价，
不能以用了神经表示就认定主题更可解释。具体推断留给文本组织一章。

事实记录逐渐增加定位范围和关系类型。DocRED（\citeyear{yao2019docred}）要求综合
跨句实体与证据判断关系；MAVEN（\citeyear{nlp-intro-wang2020maven}）标注事件触发和类型，
WikiEvents（\citeyear{nlp-intro-li2021wikievents}）把文档级论元及共指表达纳入任务。
MAVEN-ERE（\citeyear{nlp-intro-wang2022mavenere}）连接事件共指、时间、因果和子事件关系，
EventRelBench（\citeyear{nlp-intro-gong2025}）进一步检验模型对这些关系的理解。
时间先后并不推出因果，多项选择式关系判断也不能替代在文档中定位全部事件与证据。
与这条线并行，Few-NERD（\citeyear{nlp-intro-ding2021}）用细粒度类型和少样本设定
改变实体识别的监督条件，而非把少样本当成一种新的输出类型。

\subsection{文本转换与语言覆盖的研究}
\label{sec:nlp-history-generation-coverage}

文本转换一直需要回答“允许改变什么”。语法纠错可以只修局部错误，也可以为了流畅重写
整句；数据到文本要覆盖给定事实；摘要允许压缩，却不能随意改写日期和数量；翻译改变语言，
但要保持意义、术语和篇章关系。模板、统计机器翻译、神经条件生成与预训练模型改变了实现
接口，没有取消这些任务约束。

模板生成可把\code{venue=海棠楼201}填入“地点为\{venue\}”，表达选择少，但字段来源清楚。
统计机器翻译先利用词对齐估计对应，再选择短语并重排成句
（短语路线见\citeyear{nlp-intro-koehn2003}）；神经条件生成改为逐词元计算目标序列概率，
预训练模型又可在同一接口下按指令和示例适应任务。实现改变后，“4月18日”到
“April 18”的对应仍要成立，“报名截止”也不能变成“活动开始”。

内容选择、复制与编辑各有不同要求。JFLEG（\citeyear{nlp-intro-napoles2017}）
用多个人工修订研究纠错与流畅性；WebNLG（\citeyear{nlp-intro-gardent2017}）要求把给定
三元组组织成文本；XSum（\citeyear{nlp-intro-narayan2018}）强调单句高度压缩的新闻摘要；
ASSET（\citeyear{nlp-intro-alvamanchego2020}）允许替换、拆句、重排等多种简化操作。
例如“校外来宾须在4月16日18:00前提交报名信息”可简化为
“校外来宾请在4月16日18:00前报名”，但不能删成“校外来宾可直接参加”。
摘要可把旧安排压缩为“已改期”，复制机制可帮助保留地点和数字
（指针生成摘要见\citeyear{nlp-intro-see2017}），仍不能保证选择了正确的新日期。

翻译还同时扩展输出时机、语言覆盖与篇章条件。
STACL（\citeyear{nlp-intro-ma2019stacl}）在源句尚未结束时输出目标前缀，
例如听到“活动原定于4月17日”便立即给出最终日期，会在听到“现改至……”后遇到修正问题。
等待与提前输出必须共同评价，不能只按最终译文打分。
MQM式人工分析（\citeyear{nlp-intro-freitag2021}）区分准确性、流畅度、术语和上下文错误；
NLLB（\citeyear{nlp-intro-nllb2022}，预印本）扩大低资源语言的数据与评价覆盖；
WMT24（\citeyear{nlp-intro-kocmi2024}）仍在共同语言对和领域中比较专用系统、语言模型
与在线服务。DiscoX（\citeyear{nlp-intro-zhao2026discox}）把专业术语和篇章连贯纳入
中英翻译评价。把“机器学习读书会”在两段中译成同一活动名称，既需要术语约束，也需要
篇章上下文；句句流畅仍可能破坏这种一致性。语言覆盖和实时性因此不是同一种改进。

\subsection{证据使用与多轮交流的研究}
\label{sec:nlp-history-evidence-dialogue}

给定段落的抽取式问答直接找答案跨度；开放域问答还须先从资料集合找到可用段落。
SQuAD 2.0（\citeyear{rajpurkar2018squad2}）加入段落不能回答的问题，使“无答案”
成为输出规范的一部分；HotpotQA（\citeyear{yang2018hotpotqa}）要求综合多文档与支持句，
FEVER（\citeyear{nlp-intro-thorne2018}）把主张判为支持、反驳或信息不足并寻找证据。
例如“活动仍在4月17日举行”被通知反驳，“主讲人来自某学院”则是信息不足；
两种状态不能都叫“错误答案”。

\term{检索增强生成}（retrieval-augmented generation, RAG）让生成器显式读取外部文档，
为更新知识提供接口（\citeyear{lewis2020rag}）。它不同于只返回答案跨度，
可以把多个片段组织成新表述，但检索结果仍可能过时、不相关或冲突，生成器也可能
忽略证据。检索资料、组织答案与说明来源的步骤必须分别核验，RAG不是事实正确性的保证。

对话的实现也从单轮意图和槽位提取扩展到持续状态。例如“查询时间”是意图，
“ML-0418”是活动槽位；用户接着说“我是校外的”，系统应更新身份而不是丢弃此前活动。
MultiWOZ（\citeyear{budzianowski2018multiwoz}）把多领域用户约束、状态、行为和回复
组织成任务型对话资料；SGD（\citeyear{rastogi2020sgd}）用服务描述定义意图与槽位，
并检验未见服务。状态追踪决定系统当前知道什么，条件回复决定如何表达，二者并不等同。

开放交流的条件又有所不同。Persona-Chat（\citeyear{nlp-intro-zhang2018persona}）
引入说话者资料，Wizard of Wikipedia（\citeyear{nlp-intro-dinan2019}）把回复与外部知识
联系起来；人设一致不等于事实正确，引用知识也不等于已经解决用户的问题。
$\tau^2$-Bench（\citeyear{nlp-intro-barres2025}，预印本）考察用户也能改变环境时的
双边协作，$\tau$-Knowledge（\citeyear{nlp-intro-shi2026}，预印本）增加非结构化知识使用。
若用户纠正“不是帮我报名，是核对我有没有报过”，系统必须更新目标并停止写入；
工具调用成功不能抵消误解用户目标的失败。

\subsection{可检验产物与实时语音的研究}
\label{sec:nlp-history-execution-speech}

语言输出成为程序或动作后，部分正确性可以由外部过程检验。
GSM8K（\citeyear{nlp-intro-cobbe2021}，预印本）把多步应用题与候选验证联系起来；
HumanEval（\citeyear{nlp-intro-chen2021}，预印本）用执行测试评价从文档字符串生成的程序。
对“从4月16日18:00到4月18日14:00有多久”，单次生成可能只给“44小时”，
计算器或解释器还能提供$24+20=44$的可复查计算。但计算正确依赖输入时区和时间差题意，
不能把这个例子误作报名宽限期。
Tree of Thoughts（\citeyear{nlp-intro-yao2023tot}）通过候选状态、评价和搜索处理分支；
自然语言评价仍可能出错，只有明确的算术或测试约束能直接检验相应条件。

单次工具调用只解决一步交互。ReAct（\citeyear{yao2023react}）把行动与环境观察交替组织，
允许根据反馈调整后续动作；WebArena（\citeyear{nlp-intro-zhou2024webarena}）
把目标放到可复现网页环境，SWE-bench（\citeyear{nlp-intro-jimenez2024}）
要求根据问题描述修改已有代码库。后二者正式论文均为2024年，预印本均在2023年。
与生成一段貌似合理的补丁相比，测试参与的循环会记录候选、失败信息、修改及再测试；
与生成“已更新”相比，环境任务要读取更新后的状态。测试覆盖仍有边界，过程文字不是证明。

语音任务在语言内容之外增加声学与时间条件。传统识别组合声学模型、发音词典和语言模型；
CTC把无对齐音频映射为帧级路径再合并成文本，注意力识别按输出历史选择声学信息，
RNN-T联合声学时间位置与输出历史组织对齐。它们都交付转写，但允许看到多少未来音频、
什么时候可以稳定输出，并不相同，不能由模型名称直接认定支持流式识别。

合成任务则沿相反方向工作。Tacotron 2（\citeyear{nlp-intro-shen2018}）从文本预测
梅尔频谱，再用声码器生成波形；“四月十八日下午两点”读得自然仍可能把数字读错，
因此文本规范化与波形质量要分别验收。CoVoST 2（\citeyear{nlp-intro-wang2021covost}）
提供多语言声音到译文的数据条件，支持比较“识别后翻译”的级联与直接语音翻译。
Whisper（\citeyear{radford2023whisper}）研究大规模弱监督下的多语识别与翻译，
SeamlessM4T（\citeyear{nlp-intro-seamless2023}，预印本）连接多语语音和文本输入输出。

Moshi（\citeyear{defossez2024moshi}，所引为预印本）把实时双向口语交流作为对象，
$\tau$-Voice（\citeyear{nlp-intro-ray2026}，预印本）再把全双工交流放到领域任务中检验。
当用户在回复播放中说“等等，我是校外的”，系统除了识别这句话，还须处理打断、
更新状态并停止已经不合适的回复。转写错误率、译文质量、声音自然度、响应延迟和任务成功
因此必须分别评价。2026年材料在这里只提供近期问题入口，不用预印本的自述推出普遍保证。

\section*{文献说明}
\label{sec:nlp-intro-literature-notes}

以下按本章论述顺序列出评价与历史来源。它们支持相邻段落中的任务设定和接口对照，
不用于宣称领域已被穷尽，亦不为教学构造的通知、分数与轨迹背书。
已正式发表的引用保留所选出版信息，其余明确记为预印本。

\begin{fullwidth}
\noindent\term{共同评价}\par
\fullcite{kiela2021dynabench}\par
\fullcite{nlp-intro-zheng2023judge}\par
\fullcite{nlp-intro-liang2023helm}\par
\fullcite{hu2020xtreme}\par
\fullcite{bandarkar2024belebele}\par
\fullcite{singh2025globalmmlu}\par

\noindent\term{结构、判定与事实记录}\par
\fullcite{manning2008ir}\par
\fullcite{pgm-energy-models-sutton2012}\par
\fullcite{pgm-static-directed-hofmann1999}\par
\fullcite{blei2003}\par
\fullcite{nlp-intro-bianchi2021}\par
\fullcite{lee2017coreference}\par
\fullcite{nlp-intro-kitaev2018}\par
\fullcite{nlp-intro-he2018}\par
\fullcite{nivre2020universal}\par
\fullcite{nlp-intro-cer2017}\par
\fullcite{nlp-intro-williams2018}\par
\fullcite{nlp-intro-zhang2019paws}\par
\fullcite{nlp-intro-nie2020}\par
\fullcite{nlp-intro-xu2020clue}\par
\fullcite{yao2019docred}\par
\fullcite{nlp-intro-wang2020maven}\par
\fullcite{nlp-intro-li2021wikievents}\par
\fullcite{nlp-intro-wang2022mavenere}\par
\fullcite{nlp-intro-gong2025}\par
\fullcite{nlp-intro-ding2021}\par

\noindent\term{文本转换与语言覆盖}\par
\fullcite{nlp-intro-koehn2003}\par
\fullcite{nlp-intro-napoles2017}\par
\fullcite{nlp-intro-gardent2017}\par
\fullcite{nlp-intro-narayan2018}\par
\fullcite{nlp-intro-alvamanchego2020}\par
\fullcite{nlp-intro-see2017}\par
\fullcite{nlp-intro-ma2019stacl}\par
\fullcite{nlp-intro-freitag2021}\par
\fullcite{nlp-intro-nllb2022}\par
\fullcite{nlp-intro-kocmi2024}\par
\fullcite{nlp-intro-zhao2026discox}\par

\noindent\term{证据使用与多轮交流}\par
\fullcite{rajpurkar2018squad2}\par
\fullcite{yang2018hotpotqa}\par
\fullcite{nlp-intro-thorne2018}\par
\fullcite{lewis2020rag}\par
\fullcite{budzianowski2018multiwoz}\par
\fullcite{nlp-intro-zhang2018persona}\par
\fullcite{nlp-intro-dinan2019}\par
\fullcite{rastogi2020sgd}\par
\fullcite{nlp-intro-barres2025}\par
\fullcite{nlp-intro-shi2026}\par

\noindent\term{可检验产物与实时语音}\par
\fullcite{nlp-intro-cobbe2021}\par
\fullcite{nlp-intro-chen2021}\par
\fullcite{nlp-intro-yao2023tot}\par
\fullcite{yao2023react}\par
\fullcite{nlp-intro-zhou2024webarena}\par
\fullcite{nlp-intro-jimenez2024}\par
\fullcite{nlp-intro-shen2018}\par
\fullcite{nlp-intro-wang2021covost}\par
\fullcite{radford2023whisper}\par
\fullcite{nlp-intro-seamless2023}\par
\fullcite{defossez2024moshi}\par
\fullcite{nlp-intro-ray2026}\par
\end{fullwidth}

\section{本章小结}
\label{sec:nlp-introduction-summary}

自然语言处理研究的对象不是抽象的“文字能力”，而是由语言材料、任务输出、可用信息和
正确性依据共同限定的问题。同一份通知可以产生类别、跨度、结构、事实记录、新文本、
程序或环境行动；输出改变以后，位置、语义、证据和完成状态的验收方式也随之改变。

模型提供从输入到候选输出的计算接口。稀疏特征、结构化模型、神经编码器、条件生成、
外部检索和工具调用可以独立或组合使用，但模型内部表示不能代替任务结果，局部模块分数也
不能代替完整应用成功。评价需要固定语言、领域、资料时点与预算，结合直接核验、替代指标、
行为反例和分条件报告，说明证据究竟覆盖了什么。

研究历史显示，许多持续存在的任务不断增加新的输入范围、输出结构与评价条件，而不是被
某一种模型依次取代。后续章节将沿文本分析、信息抽取、分类与组织、生成与转换、问答与
对话、求解与执行以及语音任务展开。每一章都从任务对象出发，再讨论实现与评价，使“模型
能生成什么”与“应用需要证明什么”保持区分。
